%==============================================================================
%  FLEO: Fuzzy Label and Emotion Orthogonalization for YOLOv12-based FER
%  Journal manuscript (IEEEtran). Compile:
%      pdflatex fleo_journal && bibtex fleo_journal && pdflatex fleo_journal x2
%
%  NOTE ON RESULTS: every numeric cell in the results tables is a placeholder
%  written as \pending. Replace it with measured numbers from your own runs
%  (run_ablation.py / evaluate.py) before submission. Do not report
%  unmeasured values.
%==============================================================================
\documentclass[journal]{IEEEtran}

\usepackage{amsmath,amssymb,amsfonts}
\usepackage{algorithmic}
\usepackage{algorithm}
\usepackage{graphicx}
\usepackage{booktabs}
\usepackage{multirow}
\usepackage{xcolor}
\usepackage{url}
\usepackage[hidelinks]{hyperref}

% Placeholder marker for numbers that must come from real runs.
\newcommand{\pending}{\textcolor{red}{$\bullet$}}

\begin{document}

\title{FLEO: Fuzzy Label and Emotion Orthogonalization for\\
Disambiguating Morphologically Similar Expressions in\\
YOLOv12-Based Facial Emotion Recognition}

\author{Olfa~Askri,~Saifeddine~Massoud,~and~Mohamed~Ali~Hajjaji%
\thanks{O. Askri is with ENET'Com, University of Sfax, Tunisia, and with the
E$\mu$E Laboratory LR99ES30, University of Monastir, Tunisia
(e-mail: askriolfa15@gmail.com).}%
\thanks{S. Massoud and M. A. Hajjaji are with the E$\mu$E Laboratory LR99ES30,
University of Monastir, Tunisia.}%
\thanks{Manuscript submitted \today.}}

\markboth{IEEE Transactions (Submitted Manuscript)}%
{Askri \MakeLowercase{\textit{et al.}}: FLEO for YOLOv12-based Facial Emotion Recognition}

\maketitle

\begin{abstract}
Expressions that share facial musculature are routinely confused by
detectors: anger and sadness both knot the brow, fear and surprise both widen
the eyes. Inside a YOLO detector the neck aggregates channels linearly, so the
latent directions of these neighbouring classes overlap and the classification
head inherits the ambiguity. We treat this as a problem of feature geometry
rather than one to be cleaned up after the fact. FLEO is a drop-in neck module
that splits the feature tensor into $K$ per-emotion subspaces and forces them
apart with a differentiable, batched Gram--Schmidt operator, so decorrelation
is a property of the forward pass instead of a penalty negotiated with the
loss. Two design choices give the orthogonalization meaning. A lightweight
per-subspace supervision binds subspace $k$ to emotion $k$, which keeps
Gram--Schmidt from separating arbitrary channel groups; and a confusion-aware
fuzzy loss routes its smoothing mass along a data-driven confusion prior, so
the model is tolerant exactly where expressions are genuinely hard to tell
apart and strict everywhere else. We insert FLEO at the P3 and P4 neck levels
of YOLOv12 and give a deployment analysis for the Xilinx ZCU104 DPU, including
a train-time-only variant that keeps the square-root and division off the
inference path. We release a complete, reproducible implementation---model,
losses, training, a component-wise ablation harness, and the FPGA export
pipeline---and report results on FER2013 and RAF-DB. Numbers reported here are
produced by the released code; cells still marked are pending the
corresponding runs.
\end{abstract}

\begin{IEEEkeywords}
Facial emotion recognition, YOLOv12, feature orthogonalization,
Gram--Schmidt, fuzzy label smoothing, object detection, FPGA deployment.
\end{IEEEkeywords}

%==============================================================================
\section{Introduction}
%==============================================================================
\IEEEPARstart{F}{acial} emotion recognition (FER) sits between two communities
that rarely share code. One treats it as image classification on a cropped,
aligned face; the other folds detection and recognition into a single forward
pass and reports it in frames per second. The second framing is the one that
real e-learning and human--computer interaction systems need, because faces
arrive un-cropped, off-axis, and several at a time. Single-shot detectors of
the YOLO family~\cite{redmon2016yolo,khanam2024yolov11,tian2025yolov12} have
become the natural backbone for that setting, and recent work has attached FER
heads to YOLOv5 through YOLOv8 with respectable speed--accuracy
trade-offs~\cite{alshammari2024yolov8fer,ma2025alfyolo,ma2024feryolomamba}.

The accuracy ceiling in this setting is not where most papers look. It is not
the backbone, and it is not the head; it is the geometry of the neck. YOLOv12
builds its neck from area-attention and R-ELAN blocks whose output is, in the
end, a linear aggregation of channels~\cite{tian2025yolov12}. Linear
aggregation is the right tool for classes that occupy distinct regions of
feature space and the wrong one for classes that overlap there by
construction. Anger and sadness recruit overlapping brow dynamics; fear and
surprise recruit overlapping eye-widening. Their latent descriptors sit close
together, and a head that reads a linear combination of those descriptors
cannot reliably pull them apart. The FER literature has named this from the
label side---annotation ambiguity, inter-class similarity---and answered it
with uncertainty estimation and latent-distribution
mining~\cite{wang2020scn,she2021dmue,zhang2021rul,navigating2025}. Much less
has been done on the feature-geometry side inside a real-time detector.

Orthogonal decomposition is the natural instrument. Split the features into
per-class subspaces, make the subspaces orthogonal, and the head reads from
directions that no longer leak into one another. The idea is not new in
representation learning~\cite{bansal2018orthogonality,brock2019biggan,huang2018own},
and a recent transformer for near-infrared FER decomposes features into
orthogonal modality subspaces~\cite{nir2024hypergraph}. That work also makes
the observation we build on: previous orthogonal-decomposition methods impose
orthogonality through an external regularization term, which is slow to
optimize and competes with the rest of the network. An orthogonality that the
layer simply \emph{has}, rather than one the loss must negotiate, should behave
better. That is the gap FLEO targets.

FLEO---Fuzzy Label and Emotion Orthogonalization---is a single neck module
with three moving parts. It projects the neck tensor into $K$ emotion
subspaces and orthogonalizes them with a differentiable Gram--Schmidt pass, so
the constraint is structural. It binds each subspace to a specific emotion
with a small auxiliary supervision, so the orthogonalization separates
meaningful directions rather than arbitrary channel partitions; without that
binding, Gram--Schmidt on a random split is geometry for its own sake. And it
trains under a confusion-aware fuzzy loss that spreads label mass along the
actual confusion structure of the data, so the model relaxes precisely on the
pairs it has earned the right to be unsure about.

\paragraph*{Contributions}
\begin{enumerate}
\item A neck module for YOLOv12 that decomposes features into per-emotion
subspaces and enforces mutual orthogonality with a batched, differentiable
Gram--Schmidt operator, internalizing a constraint that prior work pays for
with an external penalty (Sec.~\ref{sec:method}).
\item A subspace--emotion binding loss that ties subspace $k$ to emotion $k$,
so the orthogonalization is semantically anchored rather than decorative
(Sec.~\ref{sec:binding}).
\item A confusion-aware adaptive fuzzy loss whose smoothing mass follows a
data-driven confusion prior; we state plainly its relationship to
label-smoothed cross-entropy, so the novelty claim is the structured target,
not the divergence (Sec.~\ref{sec:fuzzy}).
\item A deployment analysis for the Xilinx ZCU104 DPU that quantifies the cost
of the orthogonalization recurrence and offers a train-time-only variant with
no square-root or division on the inference path (Sec.~\ref{sec:fpga}).
\item A reproducible protocol and open implementation on FER2013 and RAF-DB
with component-wise ablations (Sec.~\ref{sec:exp}).
\end{enumerate}

%==============================================================================
\section{Related Work}
%==============================================================================
\subsection{Detection-based FER}
Treating FER as detection rather than classification buys speed and the
ability to handle multiple un-cropped faces. Alshammari and
Alshammari~\cite{alshammari2024yolov8fer} adapt YOLOv8 to seven-class
expression detection; ALF-YOLO~\cite{ma2025alfyolo} modifies YOLOv8n with
adaptive multi-branch convolutions and a large-kernel attention block;
FER-YOLO-Mamba~\cite{ma2024feryolomamba} swaps part of the attention stack for
a selective state-space model. The common thread is a stronger encoder. The
failure mode they leave untouched is the linear mixing at the neck, which is
exactly where adjacent classes blur, and none of them, to our reading, targets
YOLOv12 or the subspace geometry of the neck.

\subsection{Feature Disentanglement and Orthogonality}
Disentangling expression from identity and illumination has a long line, from
adversarial separation~\cite{ali2019degan} to orthogonality as an explicit
training objective~\cite{bansal2018orthogonality,brock2019biggan} and
orthogonal weight normalization on the Stiefel manifold~\cite{huang2018own}.
The near-infrared FER transformer of~\cite{nir2024hypergraph} decomposes
features into orthogonal modality subspaces and notes the cost of external
orthogonality penalties. FLEO differs on two axes: the orthogonalization is
computed in the forward pass, with no penalty to tune, and the subspaces are
per-emotion and explicitly bound to labels rather than per-modality.

\subsection{Label Ambiguity in FER}
Because expressions overlap and annotators disagree, one-hot targets are a
poor fit. SCN reweights uncertain samples and relabels suspected
noise~\cite{wang2020scn}; DMUE mines a latent label
distribution~\cite{she2021dmue}; RUL learns from relative sample
difficulty~\cite{zhang2021rul}; and a recent line recasts FER as partial-label
learning to keep the ambiguity rather than erase it~\cite{wang2022rethinking}.
Label smoothing is the cheap baseline~\cite{szegedy2016rethinking,muller2019labelsmoothing},
and a recent method routes smoothing toward minority
classes~\cite{navigating2025}. FLEO's loss belongs to this family but is, to
our knowledge, the first to make the smoothing target both \emph{adaptive} (a
function of per-sample uncertainty) and \emph{confusion-structured} (mass flows
along a measured confusion prior, not uniformly).

%==============================================================================
\section{Method: The FLEO Block}\label{sec:method}
%==============================================================================
FLEO is inserted at one or more neck levels of YOLOv12, after an
R-ELAN/area-attention stage and before the corresponding detection head. It
takes the level's feature tensor $X \in \mathbb{R}^{B\times C\times H\times W}$
and returns a tensor of the same shape, so it is a true drop-in.
Figure~\ref{fig:block} sketches the data path.

\begin{figure}[t]
\centering
\fbox{\parbox{0.92\columnwidth}{\centering\footnotesize
$X$ $\rightarrow$ Conv$_{1\times1}{\to}K{\cdot}d$ $\rightarrow$ Split
$\{v_1,\dots,v_K\}$ $\rightarrow$ batched Gram--Schmidt $\rightarrow$
$X_\text{ortho}$ \\[2pt]
$\odot$ SE gate $s=\sigma(\text{Conv}_{1\times1}(\text{AvgPool}(X')))$
$\rightarrow$ Conv$_{3\times3}$ $+\,X$ $\rightarrow$ $Y$ \\[2pt]
binding head $z_k=\text{GAP}(\text{Conv}_{1\times1}(e_k)) \rightarrow
\mathcal{L}_\text{bind}$
}}
\caption{The FLEO block inside the YOLOv12 neck. A $1\times1$ projection
produces $X'$, split into $K$ emotion subspaces. Batched Gram--Schmidt returns
orthonormal $e_k$. A modified squeeze-and-excitation branch supplies a fuzzy
gate $s$; the binding head taps $e_k$ for auxiliary per-subspace supervision.
A $3\times3$ convolution and a residual connection restore detector stability.}
\label{fig:block}
\end{figure}

\subsection{Decomposition into emotion subspaces}
A $1\times1$ convolution adjusts the channel count to $C_\text{mid}=K\cdot d$,
where $K$ is the number of target emotions ($K=7$ for the basic set) and
$d=\lfloor C/K\rfloor$ is the per-emotion subspace width:
\begin{equation}
X' = \text{Conv}_{1\times1}(X)\in\mathbb{R}^{B\times(K\cdot d)\times H\times W}.
\end{equation}
We then partition $X'$ along the channel axis into $K$ sub-tensors
$\mathcal{V}=\{v_1,\dots,v_K\}$, one per emotion:
\begin{equation}
v_k = X'[:,(k{-}1)d:kd,:,:]\in\mathbb{R}^{B\times d\times H\times W}.
\end{equation}
At this point the split is arbitrary; nothing yet ties $v_k$ to emotion $k$.
The binding loss of Sec.~\ref{sec:binding} does that. We flag this explicitly
because it is the step where a careless design would orthogonalize noise.

\subsection{Differentiable Gram--Schmidt orthogonalization}
To remove the shared component between adjacent emotions, we make the $K$
subspaces mutually orthogonal. For backpropagation we flatten the spatial and
channel axes of each subspace, per batch element, into a vector
$\tilde{v}_k\in\mathbb{R}^{D}$ with $D=d\cdot H\cdot W$. Classical
Gram--Schmidt then proceeds:
\begin{align}
u_1 &= \tilde{v}_1, \\
u_k &= \tilde{v}_k - \sum_{j=1}^{k-1}\text{proj}_{u_j}(\tilde{v}_k),
       \quad k=2,\dots,K, \\
\text{proj}_{u_j}(\tilde{v}_k) &= \frac{\langle\tilde{v}_k,u_j\rangle}
       {\lVert u_j\rVert_2^2+\varepsilon}\,u_j, \\
e_k &= \frac{u_k}{\lVert u_k\rVert_2+\varepsilon},
\end{align}
with $\varepsilon=10^{-8}$ for numerical safety. Each $e_k$ is reshaped back to
$\mathbb{R}^{B\times d\times H\times W}$. At the output of this stage the
property
\begin{equation}\label{eq:ortho}
\langle e_i,e_j\rangle = 0 \quad \forall\, i\neq j
\end{equation}
holds exactly per batch element. We are deliberate about what
Eq.~\eqref{eq:ortho} does and does not promise: it is a property of the
intermediate representation $X_\text{ortho}$, not of the block output $Y$. The
downstream $3\times3$ convolution and the residual re-introduce cross-subspace
mixing on purpose---the orthogonal stage acts as a structural bottleneck that
the rest of the detector is free to recombine, much as a whitening layer
conditions activations without dictating the final mapping. Treating
Eq.~\eqref{eq:ortho} as a guarantee on $Y$ would be wrong, and we do not.

A useful corollary of orthonormality is a fixed inter-class margin: for unit
vectors with $\langle e_i,e_j\rangle=0$, $\lVert e_i-e_j\rVert_2^2 =
\lVert e_i\rVert^2+\lVert e_j\rVert^2-2\langle e_i,e_j\rangle = 2$, so every
pair of emotion directions sits at distance $\sqrt{2}$ regardless of how
correlated the raw descriptors were. Our unit tests verify both
Eq.~\eqref{eq:ortho} (mean off-diagonal inner product $<10^{-10}$) and this
margin to numerical precision.

\paragraph*{Classical vs.\ modified Gram--Schmidt}
Classical Gram--Schmidt is the readable form and the less stable one. For
$K\le 8$ the conditioning is benign and the two variants agree to within
$\varepsilon$. For larger $K$, or if instability appears as exploding
gradients through the normalization, the modified variant---orthogonalize
against each $u_j$ sequentially as it is produced---is a one-line change. A
Householder or Cayley parameterization removes the division entirely; we return
to that for the FPGA in Sec.~\ref{sec:fpga}.

\subsection{Batched implementation}
The per-sample, per-$k$, per-$j$ triple loop is $\mathcal{O}(BK^2D)$ and, in
naive Python, serializes badly on a GPU. The orthogonalization vectorizes
cleanly over the batch: stack $\{\tilde{v}_k\}$ into
$V\in\mathbb{R}^{B\times K\times D}$ and run the recurrence over $k$ with
tensor operations, keeping a single Python loop of length $K$
(Alg.~\ref{alg:fleo}). The recurrence over $k$ is inherently
sequential---$u_k$ depends on $u_{1:k-1}$---but $K$ is small, so latency is set
by $K$ matrix--vector products, not by $B$.

\begin{algorithm}[t]
\caption{Batched FLEO forward pass}\label{alg:fleo}
\begin{algorithmic}[1]
\REQUIRE $X\in\mathbb{R}^{B\times C\times H\times W}$; emotions $K$;
$\varepsilon=10^{-8}$
\ENSURE $Y\in\mathbb{R}^{B\times C\times H\times W}$; binding logits $z$
\STATE $X' \leftarrow \text{Conv}_{1\times1}(X)$
\STATE $V \leftarrow \text{reshape}(\text{Split}(X',K),[B,K,D])$
       \hfill $D=dHW$
\STATE $U[:,0,:] \leftarrow V[:,0,:]$
\FOR{$k=1$ to $K-1$}
  \STATE $P \leftarrow \dfrac{V[:,k,:]\cdot U[:,:k,:]}
         {\lVert U[:,:k,:]\rVert_2^2+\varepsilon}$ \hfill coeffs, $B\times k$
  \STATE $U[:,k,:] \leftarrow V[:,k,:]-\sum_{j<k}P_j\,U[:,j,:]$
\ENDFOR
\STATE $E \leftarrow U/(\lVert U\rVert_2+\varepsilon)$ \hfill row-wise normalize
\STATE $X_\text{ortho} \leftarrow \text{reshape}(E,[B,K{\cdot}d,H,W])$
\STATE $s \leftarrow \sigma(\text{Conv}_{1\times1}(\text{AvgPool}(X')))$
       \hfill fuzzy gate
\STATE $z \leftarrow \text{GAP}(\text{Conv}_{1\times1}(E))$
       \hfill $B\times K$ binding logits
\STATE $Y \leftarrow \text{Conv}_{3\times3}(X_\text{ortho}\odot s)+X$
\RETURN $Y, z$
\end{algorithmic}
\end{algorithm}

\subsection{Fuzzy soft-assignment gate}
After orthogonalization the subspaces are recombined,
$X_\text{ortho}=\text{Concat}(e_1,\dots,e_K)$, and re-weighted by a global
channel descriptor that models partial membership across emotions---fuzzy in
the soft-assignment sense:
\begin{equation}\label{eq:gate}
s = \sigma\!\left(\text{Conv}_{1\times1}(\text{AvgPool}(X'))\right)
    \in\mathbb{R}^{B\times C_\text{mid}\times1\times1}.
\end{equation}
We compute the gate from $X'$ rather than from $X_\text{ortho}$: the gate
should read the rich pre-orthogonalization statistics, while the orthogonal
representation carries the separated directions. The block output adds a
residual to keep YOLOv12 training stable:
\begin{equation}
Y = \text{Conv}_{3\times3}(X_\text{ortho}\odot s)+X,
\end{equation}
with $\odot$ the Hadamard product broadcast over space.

\subsection{Subspace--emotion binding}\label{sec:binding}
Orthogonality is empty without semantics. If subspace $k$ is not pushed to
encode emotion $k$, Gram--Schmidt merely decorrelates seven groups of channels
that mean nothing in particular. We anchor the subspaces with a cheap
auxiliary head: each $e_k$ is reduced to a scalar by a $1\times1$ convolution
and global average pooling, giving a per-subspace logit
\begin{equation}
z_{i,k}=\text{GAP}\!\left(\text{Conv}_{1\times1}(e_{i,k})\right),
\quad z_i\in\mathbb{R}^{K},
\end{equation}
supervised against the image label $y_i$ with cross-entropy on
$\text{softmax}(z_i)$:
\begin{equation}\label{eq:bind}
\mathcal{L}_\text{bind} = -\frac{1}{B}\sum_{i=1}^{B}
\log\frac{\exp(z_{i,y_i})}{\sum_m\exp(z_{i,m})}.
\end{equation}
The head is discarded at inference. Its only job is to make the $k$-th
orthogonal direction predictive of the $k$-th emotion during training, which
is what turns Eq.~\eqref{eq:ortho} from a numerical curiosity into class
separation. The ablation in Sec.~\ref{sec:exp} is built to test exactly this:
FLEO with the binding head off should behave close to a plain SE block.

%==============================================================================
\section{Confusion-Aware Adaptive Fuzzy Loss}\label{sec:fuzzy}
%==============================================================================
Hard one-hot targets punish the model for the very confusions the geometry is
trying to model. We replace the classification term with a fuzzy target
$\hat{T}$ and a Kullback--Leibler objective:
\begin{equation}\label{eq:fuzzy}
\mathcal{L}_\text{Fuzzy}(P,\hat{T}) = \frac{1}{B}\sum_{i=1}^{B}\sum_{k=1}^{K}
\hat{T}_{i,k}\log\frac{\hat{T}_{i,k}}{P_{i,k}},
\end{equation}
with $P_{i,k}=\text{Softmax}(\text{logits}_{i,k})$.

\paragraph*{What is and is not novel}
We state the relationship to label smoothing rather than hide it. Since
$\mathrm{KL}(\hat{T}\,\|\,P)=\mathrm{CE}(\hat{T},P)-\mathcal{H}(\hat{T})$ and
the target entropy $\mathcal{H}(\hat{T})$ does not depend on the network
parameters, the gradient of Eq.~\eqref{eq:fuzzy} equals that of cross-entropy
against $\hat{T}$. With a uniform, fixed $\hat{T}$ this would be ordinary
label-smoothed cross-entropy and nothing more. The contribution is the
\emph{structure} of $\hat{T}$: it is per-sample adaptive and routed along a
confusion prior.

\paragraph*{Confusion-structured target}
Let $C\in\mathbb{R}^{K\times K}$ be a confusion prior with $C_{k,k}=0$ and
$C_{k,m}\ge 0$ measuring how confusable emotion $m$ is with $k$. We seed $C$
from an action-unit overlap table (anger--sadness and fear--surprise carry the
largest off-diagonal mass) and refresh it during training from a running
estimate of the model's own mistakes. Smoothing mass is spread over confusers
in proportion to $C$ rather than uniformly:
\begin{equation}\label{eq:target}
\hat{T}_{i,k}=
\begin{cases}
1-\alpha_i, & k=y_i,\\[4pt]
\alpha_i\,\dfrac{C_{y_i,k}}{\sum_{m\neq y_i}C_{y_i,m}}, & k\neq y_i.
\end{cases}
\end{equation}
So an angry face leaks probability toward sad, not toward happy.

\paragraph*{Adaptive ambiguity coefficient}
The smoothing strength tracks per-sample uncertainty through the prediction
entropy:
\begin{equation}\label{eq:alpha}
\alpha_i = \text{clip}\!\left(\alpha_0\,\frac{\mathcal{H}(P_i)}{\log K},
\alpha_\text{min},\alpha_\text{max}\right),
\end{equation}
with $\mathcal{H}(P_i)=-\sum_k P_{i,k}\log P_{i,k}$ and a typical range
$[\alpha_\text{min},\alpha_\text{max}]$. Confident, easy samples are trained
nearly hard; ambiguous ones are allowed to hedge along their confusers. The
target is built under a stop-gradient on $\alpha_i$, consistent with the
gradient identity above.

\paragraph*{Total objective}
In the detection setting FLEO trains end-to-end with the standard YOLOv12
localization and objectness terms plus the two FLEO terms; in the
classification setting we adopt here for the basic-emotion benchmarks, the
localization terms are inactive and the objective reduces to
\begin{equation}\label{eq:total}
\mathcal{L} = \mathcal{L}_\text{Fuzzy} + \lambda\,\mathcal{L}_\text{bind},
\end{equation}
with $\lambda$ set on validation. There is no orthogonality penalty term:
orthogonality is produced by Alg.~\ref{alg:fleo}, not purchased from the loss.

%==============================================================================
\section{Complexity and FPGA Deployment}\label{sec:fpga}
%==============================================================================
The thesis context for this work is embedded deployment on a Xilinx ZCU104
with a B4096 DPU under INT8 quantization, so we treat deployment as a
first-class constraint rather than an afterthought.

\subsection{Compute and parameter cost}
The orthogonalization adds no learnable parameters; its cost is
$\mathcal{O}(BK^2D)$ dot products and $K$ normalizations per FLEO site. The
learnable overhead is the two $1\times1$ convolutions, the SE branch, and one
$3\times3$ convolution---comparable to inserting a single SE-augmented
convolutional block. At one or two neck sites this is a small fraction of the
YOLOv12 budget. Table~\ref{tab:cost} is to be completed from profiling.

\begin{table}[t]
\caption{Cost of FLEO relative to the YOLOv12 baseline (to be filled from
profiling).}
\label{tab:cost}
\centering
\begin{tabular}{lcccc}
\toprule
Model & Params (M) & GFLOPs & GPU ms & DPU ms \\
\midrule
YOLOv12-S (baseline) & \pending & \pending & \pending & \pending \\
\quad + FLEO (P4 only) & \pending & \pending & \pending & \pending \\
\quad + FLEO (P3{+}P4) & \pending & \pending & \pending & \pending \\
\bottomrule
\end{tabular}
\end{table}

\subsection{The quantization problem, stated honestly}
Gram--Schmidt is hostile to a DPU in three specific ways, and pretending
otherwise would mislead. The normalization needs a reciprocal square root; the
projection needs a division; and the recurrence over $k$ is a sequential data
dependency. None of these maps to the DPU's INT8 convolution-centric
instruction set, so a literal port would fall back to the processing system
and stall the pipeline. We see three honest routes.

\emph{(1) Train-time-only FLEO.} Use the block during training as a regularizer
that shapes the upstream weights, then fold it out for inference. The deployed
graph keeps the projection convolutions and the SE gate---both
DPU-native---while the Gram--Schmidt pass and its square-root/division are
dropped. The hypothesis to test is that the orthogonality learned during
training survives in the weights well enough that the inference-time
approximation keeps most of the accuracy. This is the variant we recommend
profiling first, and the one our export pipeline implements.

\emph{(2) Householder/Cayley orthogonalization.} Replace Gram--Schmidt with a
product of Householder reflections or a Cayley transform, both expressible
without explicit division and friendlier to fixed-point. More invasive; kept
as an alternative.

\emph{(3) Hybrid PS/PL.} Run the $K$ small orthogonalization steps on the ARM
processing system while the DPU handles convolutions. Simple to wire, but the
round-trip can dominate latency at high frame rates.

Quantifying the accuracy--latency trade of route~(1) against the baseline is
the single most important deployment experiment, and we mark it as such in
Sec.~\ref{sec:exp}. Our released pipeline performs the Route-1 fold-out,
calibrates with the Vitis~AI quantizer, evaluates FP32 vs.\ INT8 accuracy, and
compiles to an \texttt{.xmodel} for the DPUCZDX8G/ZCU104
target~\cite{vitisai}.

%==============================================================================
\section{Experimental Protocol and Results}\label{sec:exp}
%==============================================================================
\textit{Note.} Numeric cells marked \pending\ are placeholders to be populated
from the released code before submission; we do not report unmeasured values.

\subsection{Datasets}
We evaluate on two complementary benchmarks. \textbf{FER2013}~%
\cite{goodfellow2013fer2013}: 35{,}887 grayscale $48\times48$ images, seven
classes, noisy labels, with a human ceiling in the mid-70s. \textbf{RAF-DB}~%
\cite{li2017rafdb}: roughly 15{,}000 in-the-wild images with crowdsourced
labels on the basic-emotion split. Both are classification corpora; we cast
each labelled face as a single, full-frame target so a YOLOv12 backbone with
the FLEO neck can be trained and read out by the pooled emotion head, and we
report standard FER accuracy and macro-F1. AffectNet~\cite{mollahosseini2019affectnet}
is the natural third benchmark if compute allows.

\subsection{Implementation details}
The backbone is a YOLOv12 feature extractor tapped at the P3 ($/8$) and P4
($/16$) levels, with a FLEO block on each and a global-pooling emotion head.
Unless noted we use $K=7$, per-subspace width $d$ from the level channel count,
$\varepsilon=10^{-8}$, and input resolution $128\times128$ (grayscale
replicated to three channels). Optimization uses AdamW~\cite{loshchilov2019adamw}
with cosine annealing, weight decay $10^{-4}$, and gradient-norm clipping at
$10$. Because FER2013 is strongly imbalanced (the \emph{disgust} class is an
order of magnitude smaller than \emph{happy}), training draws from a
class-balanced sampler; augmentation is horizontal flip, mild rotation,
resized crop, colour jitter, and random erasing. The confusion prior $C$ is
seeded from action-unit overlap and updated each epoch by an exponential
moving average of the model's mistakes. During training and validation the
Gram--Schmidt pass is active in both passes so that the two match; the
train-time-only fold-out is applied only for FPGA export. We report
mean~$\pm$~std over three seeds.

\subsection{Baselines and metrics}
A YOLOv12 detector without FLEO is the primary baseline. We also compare
against an SE-only neck block, to separate the gain of orthogonality from the
gain of channel gating. Metrics are overall accuracy and macro-F1 (the
imbalance makes macro-F1 the load-bearing number), per-class precision/recall,
the seven-class confusion matrix, and---for the deployment study---INT8 vs.\
FP32 accuracy and latency on the ZCU104.

\subsection{Main results}
Table~\ref{tab:main} reports accuracy and macro-F1 on both benchmarks.

\begin{table}[t]
\caption{FER accuracy and macro-F1 (mean over three seeds). Cells marked
\pending\ are pending the corresponding runs.}
\label{tab:main}
\centering
\begin{tabular}{lcccc}
\toprule
\multirow{2}{*}{Method} & \multicolumn{2}{c}{FER2013} &
\multicolumn{2}{c}{RAF-DB}\\
\cmidrule(lr){2-3}\cmidrule(lr){4-5}
 & Acc. & m-F1 & Acc. & m-F1 \\
\midrule
YOLOv12 (baseline)          & \pending & \pending & \pending & \pending \\
\quad + SE neck             & \pending & \pending & \pending & \pending \\
\quad + FLEO (no binding)   & \pending & \pending & \pending & \pending \\
\quad + FLEO (full)         & \pending & \pending & \pending & \pending \\
\quad + FLEO (full) + fuzzy & \pending & \pending & \pending & \pending \\
\bottomrule
\end{tabular}
\end{table}

\subsection{Ablations}
The table is designed so each row isolates one claim.
\emph{(i)}~SE-only vs.\ FLEO-no-binding tests whether orthogonality helps at
all. \emph{(ii)}~FLEO-no-binding vs.\ FLEO-full tests the central thesis that
binding is what makes orthogonality meaningful; if these two rows are equal,
the binding head is dead weight and the paper should say so.
\emph{(iii)}~$\pm$ fuzzy loss isolates the loss from the architecture.
\emph{(iv)}~A confusion-matrix comparison on the anger/sadness and
fear/surprise cells, baseline vs.\ FLEO, is the qualitative evidence that the
targeted pairs separate; Grad-CAM~\cite{selvaraju2017gradcam} and a
t-SNE~\cite{maaten2008tsne} of the $\{e_k\}$ make the same point visually. The
released \texttt{run\_ablation.py} produces rows (i)--(iii) directly, and
\texttt{evaluate.py} produces the confusion matrix and t-SNE.

\subsection{Deployment experiment}
We profile the three deployment routes of Sec.~\ref{sec:fpga} on the ZCU104.
The decisive number is the accuracy drop of train-time-only FLEO (route~1)
against full FLEO, traded against its DPU latency; we report INT8 vs.\ FP32
accuracy and frames per second in Table~\ref{tab:deploy}.

\begin{table}[t]
\caption{Route-1 (train-time-only) deployment on ZCU104 (to be filled).}
\label{tab:deploy}
\centering
\begin{tabular}{lccc}
\toprule
Configuration & Acc. (\%) & FPS & DPU ms \\
\midrule
Full FLEO (FP32, GPU)       & \pending & \pending & -- \\
Route-1 FLEO (FP32)         & \pending & \pending & \pending \\
Route-1 FLEO (INT8, DPU)    & \pending & \pending & \pending \\
\bottomrule
\end{tabular}
\end{table}

%==============================================================================
\section{Limitations}
%==============================================================================
A few risks are worth naming before a reviewer does. The orthogonality holds
at $X_\text{ortho}$, not at $Y$; we argue the bottleneck still helps, but that
is an empirical claim the ablation must support, not a theorem. The fuzzy loss
reduces, in gradient, to structured label smoothing, so its benefit rests
entirely on the confusion prior $C$ being a better target than a uniform
one---if $C$ is poorly estimated, the loss is label smoothing with extra
steps. Finally, the version of the method that is cleanest for accuracy (full
Gram--Schmidt) is the hardest to deploy, while the one that is cleanest to
deploy (train-time-only) is the one whose accuracy is unproven until the
deployment experiment is run. These are the three places the work can fail, and
the protocol targets each.

%==============================================================================
\section{Conclusion}
%==============================================================================
FER inside a real-time detector fails where adjacent expressions share
musculature, and the failure is geometric: the neck mixes their latent
directions linearly. FLEO answers with a structural orthogonalization of
per-emotion subspaces, made meaningful by a binding head and made tolerant by
a confusion-aware fuzzy loss. The orthogonality is internal to the forward
pass, so no penalty needs tuning, and the deployment story is told honestly,
with a train-time-only variant for the DPU and a clear statement of what it has
to prove. The implementation, the ablation harness, and the FPGA export are
released so that the runs reported here are reproducible and the remaining
placeholders can be closed by anyone with the two public benchmarks.

\section*{Acknowledgment}
The authors thank the E$\mu$E Laboratory, University of Monastir, for
computational support. % [TODO: funding/grant acknowledgments]

\bibliographystyle{IEEEtran}
\bibliography{references}

\end{document}
